% =====================================================================
%  main.tex : 本書全体
%    ビルド：cd tex && latexmk        （出力は tex/build/main.pdf）
%    一部の章だけ作るときは、下の \includeonly のコメントを外す。
% =====================================================================
% =====================================================================
%  preamble.tex : main.tex と sample.tex で共通のプリアンブル
% =====================================================================
\documentclass[
  book,                 % 片面ではなく見開き（章は奇数ページから）
  paper=b5j,            % JIS B5（182mm × 257mm）
  fontsize=10pt,
  line_length=40zw,     % 1行 40 字
  number_of_lines=34,   % 1ページ 34 行
  gutter=22mm,          % のど（とじ側）の余白
  head_space=24mm,      % 天（上）の余白
  baselineskip=1.7zw,
]{jlreq}

\usepackage{style/textbook}


% \includeonly{parts/statistics/02-probability}

\hypersetup{
  pdftitle={統計検定 データサイエンスエキスパート 教科書},
}

\begin{document}

\frontmatter
% 扉（タイトルページ）
\begin{titlepage}
  \centering
  \vspace*{0.22\textheight}
  {\sffamily\gtfamily\bfseries\Large 統計検定\par}
  \vspace{0.8\baselineskip}
  {\sffamily\gtfamily\bfseries\Huge\color{tbMain} データサイエンスエキスパート\par}
  \vspace{0.6\baselineskip}
  {\sffamily\gtfamily\bfseries\LARGE 教科書\par}
  \vspace{1.2\baselineskip}
  {\color{tbMain}\rule{0.6\textwidth}{1pt}\par}
  \vspace{1.2\baselineskip}
  {\large 高校数学と Python から始める\par}
  \vfill
  {\large \today\par}
\end{titlepage}

\chapter{はじめに}

\section*{この本のねらい}

この本は、統計検定「データサイエンスエキスパート」の出題範囲を、
\textbf{高校数学と Python の文法だけ}を前提に学べるように書いた教科書です。
大学で学ぶ数学（行列やガンマ関数など）が必要になる所では、その場で最小限の説明を入れています。

章立ては、試験の出題範囲表に合わせています。
範囲表の「大項目」を部、「中項目」を章、「小項目」を節としました。
各節の最初には、範囲表に書かれているキーワードを示しています。

\section*{この本の読み方}

新しい用語が初めて出てくる所では、{\sffamily\gtfamily\bfseries 標本平均}（sample mean）のように太字で示し、英語の名前も添えます。
英語の名前は、英語の資料を読むときや、Python のライブラリで関数を探すときに役立ちます。
太字の用語は、巻末の索引にも載せています。

本文の中では、次のような囲みを使います。
定義・定理・例・例題には、章ごとの通し番号を付けます。
たとえば「定理 3.2」は、第 3 章の 2 番目の囲みという意味です。

\begingroup
\renewcommand{\thechapter}{3}  % 見本として「第3章」の番号を表示する

\begin{definition}
  言葉の意味を決めるものです。後の議論はすべてここから始まります。
\end{definition}

\begin{theorem}
  定義から導かれる、大事な性質です。この本では、定理を必ず\textbf{証明}（導出）します。
\end{theorem}

\begin{intuition}
  数式や定理が何を言っているのかを、言葉や図で説明します。
\end{intuition}

\begin{merit}
  その定理や数式を知っていると、何が楽になるのか、何ができるようになるのかを説明します。
  式を覚えるより先に、ここを読んでください。
\end{merit}

\begin{example}
  具体的な数値や場面で、定義や定理を確かめます。
\end{example}

\begin{problem}
  自分で手を動かして解いてみる問題です。
  \tcblower
  解答は、点線の下にあります。
\end{problem}
\endgroup

\begin{mathnote}
  高校数学を超える道具を、その場で必要なだけ説明します。
  くわしくは「数学基礎」の部で扱います。
\end{mathnote}

\begin{supplement}
  範囲表には書かれていませんが、後の内容を理解するために必要な話です。
\end{supplement}

\begin{reference}
  発展的な内容です。名前と大まかな意味を知っていれば十分です。
\end{reference}

\begin{exampoint}
  試験でよく問われる点や、まちがえやすい点です。
\end{exampoint}

Python のコードは次のように示します。
コードはすべて \texttt{numpy} などの基本的なライブラリだけで動くようにしています。

\begin{pycode}
import numpy as np
print(np.mean([1, 2, 3]))
\end{pycode}

\begin{pyoutput}
2.0
\end{pyoutput}

\section*{記号について}

\begin{center}
  \begin{tabular}{ll}
    \toprule
    記号 & 意味 \\
    \midrule
    $\Prob(A)$ & 事象 $A$ が起こる確率 \\
    $\E[X]$, $\V[X]$ & 確率変数 $X$ の期待値、分散 \\
    $\Cov(X,Y)$ & $X$ と $Y$ の共分散 \\
    $X \sim \Normal(\mu,\sigma^2)$ & $X$ が平均 $\mu$、分散 $\sigma^2$ の正規分布にしたがう \\
    $X_1,\dots,X_n \iid F$ & $X_1,\dots,X_n$ が互いに独立に、同じ分布 $F$ にしたがう \\
    $\bX$ & 標本平均 $\frac{1}{n}\sum_{i=1}^n X_i$ \\
    $\estim{\theta}$ & パラメータ $\theta$ の推定量 \\
    \bottomrule
  \end{tabular}
\end{center}

\tableofcontents

\mainmatter

% ---------------------------------------------------------------------
\part{統計基礎}
\label{part:statistics}
% =====================================================================
%  統計基礎 0. 準備：確率と統計の基本（範囲表外・2級相当）
%  構成は Docs/SectionPlan.md にしたがう。
% =====================================================================
\chapter{確率と統計の基本}
\label{chap:stat-preparation}

\begin{chaptergoal}
  % TODO: この章の目標
\end{chaptergoal}

% 範囲表には載っていない（2級相当）。範囲表の注記「2級・3級などの内容も出題される」に対応する。

% ---------------------------------------------------------------------
\section{データの要約}
\label{sec:stat-summary}

% TODO: Σ（シグマ）記号の使い方
% TODO: 平均・分散・標準偏差
% TODO: 共分散と相関係数
% TODO: 度数分布表、ヒストグラム、箱ひげ図

% ---------------------------------------------------------------------
\section{確率の基本}
\label{sec:stat-probability-basics}

% TODO: 試行と事象、和事象・積事象・余事象
% TODO: 確率の加法定理
% TODO: 条件付き確率と乗法定理
% TODO: 事象の独立
% TODO: ベイズの定理（簡単な例：検査の陽性と病気の確率）

% ---------------------------------------------------------------------
\section{確率変数と期待値}
\label{sec:stat-random-variable}

% TODO: 離散型の確率変数と確率分布
% TODO: 連続型の確率変数と確率密度関数（「面積が確率」というイメージ）
% TODO: 累積分布関数
% TODO: 期待値と分散の定義と性質（$E[aX+b]$、$V[aX+b]$ など）
% TODO: 2つの確率変数：同時分布、周辺分布、共分散、独立
% TODO: 和の期待値と分散（独立なときの分散の足し算）

% ---------------------------------------------------------------------
\section{基本的な確率分布}
\label{sec:stat-basic-distributions}

% TODO: 離散型：ベルヌーイ分布、二項分布、ポアソン分布、幾何分布
% TODO: 連続型：一様分布、指数分布、正規分布
% TODO: 標準化と標準正規分布表の使い方

% ---------------------------------------------------------------------
\section{推定と検定の基本}
\label{sec:stat-estimation-testing-basics}

% TODO: 母集団と標本、統計量
% TODO: 標本平均の期待値と分散
% TODO: 推定量と不偏性（不偏分散で $n-1$ で割る理由）
% TODO: 仮説検定の流れ：帰無仮説と対立仮説、有意水準、棄却域、p 値
% TODO: 第1種の過誤・第2種の過誤、検出力
% TODO: 片側検定と両側検定

\begin{chaptersummary}
  % TODO: この章のまとめ
\end{chaptersummary}

% =====================================================================
%  統計基礎 1. 確率と確率分布
%  構成は Docs/SectionPlan.md にしたがう。
% =====================================================================
\chapter{確率と確率分布}
\label{chap:stat-probability}

\begin{chaptergoal}
  % TODO: この章の目標
\end{chaptergoal}

% ---------------------------------------------------------------------
\section{確率分布、確率変数}
\label{sec:stat-moments}

\keywords{チェビシェフの不等式、積率、尖度、歪度、積率母関数}

\subsection{積率（モーメント）}
% TODO:   原点のまわりの $k$ 次積率 $E[X^k]$ と、平均のまわりの $k$ 次積率 $E[(X-\mu)^k]$
% TODO:   期待値と分散を積率で表す

\subsection{歪度（わいど）}
% TODO:   3次の積率を標準化したもの
% TODO:   分布の左右のかたよりを表す。正・負・0 の場合を図で比べる

\subsection{尖度（せんど）}
% TODO:   4次の積率を標準化したもの
% TODO:   正規分布では 3 になる。3 を引いた「超過尖度」の流儀もあることに注意
% TODO:   分布のとがり具合と、裾（すそ）の重さを表す

\subsection{積率母関数}
% TODO:   定義 $M_X(t)=E[e^{tX}]$
% TODO:   微分して $t=0$ を代入すると積率が得られる
% TODO:   積率母関数が同じなら分布も同じ（分布を見分ける道具になる）
% TODO:   独立な確率変数の和の積率母関数は、それぞれの積になる
% TODO:   例：二項分布、ポアソン分布、正規分布の積率母関数
% TODO:   〔数学メモ〕$e^x$ の微分、テイラー展開のイメージ

\subsection{チェビシェフの不等式}
% TODO:   主張：$P(|X-\mu|\ge k\sigma)\le 1/k^2$
% TODO:   証明（マルコフの不等式を経由する）
% TODO:   分布の形を知らなくても使えるという利点
% TODO:   正規分布での実際の確率と比べると、評価がゆるいこと

\subsection{確率変数の変換}
% 〔補足〕
% TODO:   $Y=g(X)$ の確率密度関数の求め方
% TODO:   対数正規分布（1.2）や逆変換法（4.2）の準備

% ---------------------------------------------------------------------
\section{主要な確率分布}
\label{sec:stat-distributions}

\keywords{対数正規分布、ガンマ分布、ベータ分布、超幾何分布、負の二項分布、多変量正規分布}

\subsection{超幾何分布}
% TODO:   「戻さずに取り出す」（非復元抽出）ときの当たりの個数の分布
% TODO:   期待値と分散
% TODO:   二項分布との関係（母集団が大きいと二項分布に近づく）、有限母集団修正

\subsection{負の二項分布}
% TODO:   「$r$ 回成功するまでに失敗した回数」の分布
% TODO:   幾何分布を一般にしたもの
% TODO:   期待値と分散
% TODO:   使われる場面：ポアソン分布よりばらつきが大きいカウントデータ（過分散）

\subsection{ガンマ分布}
% TODO:   〔数学メモ〕ガンマ関数 $\Gamma(\alpha)$：広義積分による定義、$\Gamma(n)=(n-1)!$、$\Gamma(1/2)=\sqrt{\pi}$（数学基礎「1変数関数の積分法」を参照）
% TODO:   確率密度関数、形状パラメータと尺度パラメータ
% TODO:   期待値と分散
% TODO:   指数分布の和はガンマ分布になる
% TODO:   カイ二乗分布はガンマ分布の特別な場合（2.1 への準備）

\subsection{ベータ分布}
% TODO:   〔数学メモ〕ベータ関数 $B(a,b)$ と、ガンマ関数との関係
% TODO:   0 から 1 の値をとる分布で、「確率そのもの」の分布として使える
% TODO:   一様分布は特別な場合
% TODO:   期待値と、パラメータによる形の変化
% TODO:   ベイズ統計の事前分布として使うことの予告（3.1）

\subsection{対数正規分布}
% TODO:   $\log X$ が正規分布にしたがうときの $X$ の分布
% TODO:   右に長く裾をひく形。所得や株価などの例
% TODO:   期待値 $e^{\mu+\sigma^2/2}$、中央値 $e^{\mu}$

\subsection{多変量正規分布}
% TODO:   〔数学メモ〕ベクトルと行列による表し方、分散共分散行列（数学基礎「線形代数」を参照）
% TODO:   2変量正規分布の確率密度関数と相関係数
% TODO:   等高線が楕円になること
% TODO:   周辺分布・条件付き分布も正規分布になる
% TODO:   正規分布では「無相関 ⇔ 独立」が成り立つ（一般には成り立たない）
% TODO:   線形変換しても多変量正規分布になる

\subsection{まとめ}
% ：分布どうしの関係図（二項・ポアソン・正規・ガンマ・カイ二乗・ベータ など）

% ---------------------------------------------------------------------
\section{確率変数の漸近的性質}
\label{sec:stat-asymptotics}

\keywords{大数の法則、中心極限定理、確率収束、分布収束}

\subsection{確率収束}
% TODO:   定義と直感的な意味（「ずれが大きくなる確率が 0 に近づく」）

\subsection{大数の法則}
% TODO:   主張：標本平均は母平均に確率収束する
% TODO:   チェビシェフの不等式を使った証明
% TODO:   コイン投げのシミュレーションで確かめる

\subsection{分布収束}
% TODO:   定義（累積分布関数が近づく）
% TODO:   確率収束との違い（確率収束なら分布収束、逆は成り立たない）

\subsection{中心極限定理}
% TODO:   主張：標本平均を標準化すると、標準正規分布に分布収束する
% TODO:   もとの分布が正規分布でなくてよいこと
% TODO:   二項分布の正規近似
% TODO:   積率母関数を使った証明の流れ（概略）
% TODO:   注意：分散が存在しない分布（コーシー分布）では成り立たない

\begin{chaptersummary}
  % TODO: この章のまとめ
\end{chaptersummary}

% =====================================================================
%  統計基礎 2. 推測統計
%  構成は Docs/SectionPlan.md にしたがう。
% =====================================================================
\chapter{推測統計}
\label{chap:stat-inference}

\begin{chaptergoal}
  % TODO: この章の目標
\end{chaptergoal}

% ---------------------------------------------------------------------
\section{標本分布}
\label{sec:stat-sampling-distribution}

\keywords{カイ二乗分布、標本平均と標本分散の独立性、t 分布、F 分布}

% TODO: 〔補足〕標本分布とは何か（統計量も確率変数であること）

\subsection{カイ二乗分布}
% TODO:   定義：独立な標準正規分布にしたがう変数の2乗の和
% TODO:   自由度の意味、期待値と分散、形の変化
% TODO:   ガンマ分布との関係、再生性

\subsection{標本平均と標本分散の独立性}
% TODO:   正規母集団では、標本平均 $\bar X$ と標本分散 $S^2$ が独立になる
% TODO:   $(n-1)S^2/\sigma^2$ が自由度 $n-1$ のカイ二乗分布にしたがう
% TODO:   自由度が1つ減る理由の直感（偏差の合計が必ず 0 になる）
% TODO:   証明の流れ（直交変換を使う。数学基礎「線形代数」を参照）

\subsection{t 分布}
% TODO:   定義：$Z/\sqrt{V/k}$（$Z$ は標準正規、$V$ は自由度 $k$ のカイ二乗、互いに独立）
% TODO:   母分散がわからないときの標本平均の分布として現れる
% TODO:   形の特徴（正規分布より裾が重い）、自由度を大きくすると正規分布に近づく

\subsection{F 分布}
% TODO:   定義：2つの独立なカイ二乗変数を、それぞれの自由度で割った比
% TODO:   2つの母分散の比の推測、分散分析（2.4）での使い方
% TODO:   t 分布との関係（$t^2$ は F 分布にしたがう）

% ---------------------------------------------------------------------
\section{点推定、区間推定}
\label{sec:stat-estimation}

\keywords{一致性、有効性、信頼区間と信頼係数}

\subsection{点推定の方法}
% 〔補足〕
% TODO:   モーメント法
% TODO:   最尤法：尤度関数、対数尤度、微分して 0 とおく（尤度比検定・EM 法の準備）

\subsection{推定量のよさの基準}
% TODO:   〔補足〕不偏性（復習）と平均二乗誤差（＝バイアスの2乗＋分散）
% TODO:   一致性：標本の大きさを増やすと真の値に確率収束すること。大数の法則との関係と例
% TODO:   有効性：不偏推定量の中で分散が最小であること
% TODO:     フィッシャー情報量とクラメール・ラオの下限
% TODO:     下限に達する推定量（有効推定量）の例
% TODO:   〔参考〕最尤推定量の漸近的な性質（一致性・漸近正規性）

\subsection{区間推定}
% TODO:   信頼区間と信頼係数の意味
% TODO:     「95% 信頼区間」の正しい解釈（同じ方法で何度も区間を作ると、その 95% が真の値を含む）
% TODO:     よくある誤解
% TODO:   正規母集団の母平均（母分散が既知のとき・未知のとき）
% TODO:   正規母集団の母分散（カイ二乗分布を使う）
% TODO:   母比率（正規近似）
% TODO:   2つの母平均の差、2つの母分散の比

% ---------------------------------------------------------------------
\section{汎用的な検定}
\label{sec:stat-general-tests}

\keywords{尤度比検定、ノンパラメトリック検定、ウィルコクソン検定、並べ替え検定}

\subsection{尤度比検定}
% TODO:   尤度比の定義と考え方（帰無仮説のもとでの最大尤度と、全体での最大尤度の比）
% TODO:   〔参考〕ネイマン・ピアソンの補題（最も検出力の高い検定）
% TODO:   $-2\log(\text{尤度比})$ が近似的にカイ二乗分布にしたがうこと（ウィルクスの定理）
% TODO:   例：正規母集団の母平均の検定、二項分布の比率の検定

\subsection{ノンパラメトリック検定}
% TODO:   パラメトリック検定との違い（母集団の分布を仮定しない）
% TODO:   長所（外れ値に強い）と短所（検出力が下がることがある）
% TODO:   順位を使うという考え方、符号検定

\subsection{ウィルコクソン検定}
% TODO:   順位和検定：対応のない2群の比較（マン・ホイットニーの U 検定と同じもの）
% TODO:   符号付き順位検定：対応のある2群の比較
% TODO:   小さなデータでの計算例、標本が大きいときの正規近似

\subsection{並べ替え検定}
% TODO:   考え方：帰無仮説が正しければ、群のラベルを入れかえても分布は変わらない
% TODO:   手順と、正確な p 値の求め方
% TODO:   組合せの数が多いときはランダムに入れかえて近似する（計算統計 4 章への橋渡し）

% ---------------------------------------------------------------------
\section{種々の検定}
\label{sec:stat-various-tests}

\keywords{一元配置分散分析、二元配置分散分析、交互作用、適合度検定}

\subsection{一元配置分散分析}
% TODO:   モデル：データ ＝ 全体平均 ＋ 群の効果 ＋ 誤差
% TODO:   平方和の分解（全体の平方和 ＝ 群間平方和 ＋ 群内平方和）
% TODO:   自由度、分散分析表、F 検定
% TODO:   t 検定との関係（2群のとき）

\subsection{二元配置分散分析}
% TODO:   2つの要因の主効果
% TODO:   繰り返しのない場合と、繰り返しのある場合
% TODO:   平方和の分解と分散分析表

\subsection{交互作用}
% TODO:   意味：一方の要因の効果が、もう一方の要因の水準によって変わること
% TODO:   交互作用プロットによる見方
% TODO:   繰り返しのある二元配置での交互作用の検定

\subsection{適合度検定}
% TODO:   ピアソンのカイ二乗統計量
% TODO:   自由度の決め方（推定したパラメータの数だけ減る）
% TODO:   例：サイコロの公平性、ポアソン分布への当てはまり
% TODO:   分割表の独立性の検定
% TODO:   〔参考〕尤度比検定（G 検定）との関係

% ---------------------------------------------------------------------
\section{多重比較}
\label{sec:stat-multiple-comparison}

\keywords{ボンフェロニ補正}

% TODO: 多重性の問題：検定を何回もくり返すと、どこかで誤って棄却する確率が大きくなる（$1-(1-\alpha)^m$ の計算例）
% TODO: ファミリーワイズ・エラー率（FWER）

\subsection{ボンフェロニ補正}
% TODO:   方法：有意水準を検定の回数 $m$ で割る（$\alpha/m$）
% TODO:   根拠：確率の和の不等式（ブールの不等式）
% TODO:   欠点：保守的（棄却しにくくなる）
% TODO: 〔参考〕ホルム法、テューキー法、偽発見率（FDR）

\begin{chaptersummary}
  % TODO: この章のまとめ
\end{chaptersummary}

% =====================================================================
%  統計基礎 3. ベイズ理論
%  構成は Docs/SectionPlan.md にしたがう。
% =====================================================================
\chapter{ベイズ理論}
\label{chap:stat-bayes}

\begin{chaptergoal}
  % TODO: この章の目標
\end{chaptergoal}

% ---------------------------------------------------------------------
\section{事前分布・事後分布}
\label{sec:stat-prior-posterior}

\keywords{事前分布、共役事前分布、事後分布}

% TODO: 〔補足〕ベイズ統計の考え方
% TODO:   ベイズの定理の復習
% TODO:   パラメータを確率変数として扱う。これまでの（頻度論の）考え方との違い

\subsection{事前分布}
% TODO:   データを見る前の知識や信念を分布で表す
% TODO:   情報が少ないときの事前分布（無情報事前分布）

\subsection{事後分布}
% TODO:   事後分布 ∝ 尤度 × 事前分布
% TODO:   正規化定数（全体の確率を 1 にする定数）の役割
% TODO:   データが増えると事前分布の影響が小さくなること

\subsection{共役事前分布}
% TODO:   定義：事前分布と事後分布が同じ種類の分布になる組み合わせ
% TODO:   ベータ分布と二項分布（コイン投げの例）
% TODO:   ガンマ分布とポアソン分布
% TODO:   正規分布と正規分布（分散が既知の場合）
% TODO: 〔補足〕事後分布のまとめ方：事後平均、事後確率最大の値（MAP 推定）、信用区間（信頼区間との違い）

% ---------------------------------------------------------------------
\section{ベイズ的仮説検定}
\label{sec:stat-bayes-testing}

\keywords{ベイズファクター、ベイズ判別（各カテゴリーの事後確率）}

% TODO: 考え方：仮説そのものが正しい確率（事後確率）を比べる
% TODO: 〔補足〕周辺尤度（パラメータについて平均した尤度）

\subsection{ベイズファクター}
% TODO:   定義：2つの仮説（モデル）の周辺尤度の比
% TODO:   事後オッズ ＝ ベイズファクター × 事前オッズ
% TODO:   値の大きさの目安（ジェフリーズの基準）
% TODO:   p 値との違い

\subsection{ベイズ判別}
% TODO:   各カテゴリーの事後確率を、ベイズの定理で計算する
% TODO:   事後確率が最大のカテゴリーに分類する（誤分類の確率が最小になる）
% TODO:   例：2つの正規分布のどちらから来たデータかを判別する
% TODO:   モデリング側の「単純ベイズ」「ベイズ判別」への橋渡し

\begin{chaptersummary}
  % TODO: この章のまとめ
\end{chaptersummary}

% =====================================================================
%  統計基礎 4. 計算統計
%  構成は Docs/SectionPlan.md にしたがう。
% =====================================================================
\chapter{計算統計}
\label{chap:stat-computational}

\begin{chaptergoal}
  % TODO: この章の目標
\end{chaptergoal}

% ---------------------------------------------------------------------
\section{ブートストラップ}
\label{sec:stat-bootstrap}

\keywords{復元抽出、経験分布、リサンプリング}

% TODO: 動機：推定量の分布を式で求めるのが難しい場合がある

\subsection{経験分布}
% TODO:   各データに等しい確率 $1/n$ をおいた分布
% TODO:   経験分布関数と、真の分布関数との関係

\subsection{復元抽出}
% TODO:   取り出したものを戻してから次を取り出す抽出（非復元抽出との違い）

\subsection{リサンプリング}
% TODO:   手元のデータから復元抽出をくり返して「ブートストラップ標本」を作る
% TODO:   標準誤差の推定
% TODO:   パーセンタイル法による信頼区間
% TODO:   手順の整理と、Python による簡単な実装例
% TODO: 〔参考〕ジャックナイフ法

% ---------------------------------------------------------------------
\section{サンプリング}
\label{sec:stat-sampling}

\keywords{擬似乱数、逆変換法、棄却法、マルコフ連鎖モンテカルロ法}

\subsection{擬似乱数}
% TODO:   コンピュータで作る「乱数のように見える数」
% TODO:   線形合同法のしくみ、〔参考〕メルセンヌ・ツイスタ
% TODO:   シード（種）と再現性

\subsection{逆変換法}
% TODO:   一様乱数を累積分布関数の逆関数に入れると、目的の分布の乱数になる
% TODO:   例：指数分布の乱数
% TODO:   離散分布の場合

\subsection{棄却法}
% TODO:   作りやすい分布（提案分布）から乱数を作り、一定の確率で採用する
% TODO:   採用する確率の決め方と、なぜ正しい分布になるか
% TODO:   効率（採用される割合）

\subsection{マルコフ連鎖モンテカルロ法（MCMC）}
% TODO:   〔補足〕マルコフ連鎖：次の状態が今の状態だけで決まる確率的な動き、推移確率、定常分布
% TODO:   考え方：目的の分布を定常分布にもつマルコフ連鎖を作り、長く動かして乱数を得る
% TODO:   メトロポリス・ヘイスティングス法
% TODO:   ギブスサンプリング
% TODO:   バーンイン（初めの部分を捨てる）、収束の確かめ方
% TODO:   ベイズ統計の事後分布からのサンプリングへの応用

% ---------------------------------------------------------------------
\section{モンテカルロ積分}
\label{sec:stat-monte-carlo}

\keywords{モンテカルロ積分、期待値や確率密度の正規化定数、分散減少法}

\subsection{モンテカルロ積分}
% TODO:   積分を「期待値」として表し、乱数の平均で近似する
% TODO:   例：円周率 $\pi$ の近似
% TODO:   誤差は $1/\sqrt{n}$ の速さで小さくなる（中心極限定理による）
% TODO:   次元が高い積分でも使いやすいという利点

\subsection{期待値や確率密度の正規化定数}
% TODO:   期待値の近似計算
% TODO:   確率密度関数の正規化定数（全体を 1 にする定数）を乱数で近似する
% TODO:   ベイズ統計の周辺尤度の計算との関係

\subsection{分散減少法}
% TODO:   同じ乱数の数で、近似の誤差を小さくする工夫
% TODO:   重点サンプリング（重要な所から多く乱数をとる）
% TODO:   制御変量法
% TODO:   負の相関法（対称変量法）
% TODO:   層別サンプリング

% ---------------------------------------------------------------------
\section{欠測値の扱い}
\label{sec:stat-missing}

\keywords{EM 法}

% TODO: 〔補足〕欠測値の種類
% TODO:   完全にランダムな欠測（MCAR）、ランダムな欠測（MAR）、ランダムでない欠測（MNAR）
% TODO:   欠測のある行を捨てる方法（完全ケース分析）や平均で埋める方法の問題点

\subsection{EM 法}
% TODO:   考え方：欠けている部分を「期待値で補う（E ステップ）」と「パラメータを最尤推定する（M ステップ）」をくり返す
% TODO:   例1：一部が欠測した正規分布のデータ
% TODO:   例2：混合正規分布（どのグループから来たかが分からないデータ）
% TODO:   くり返すごとに尤度が減らないこと
% TODO:   注意：局所的な最大値に止まることがある、初期値の選び方
% TODO: 〔参考〕多重代入法

\begin{chaptersummary}
  % TODO: この章のまとめ
\end{chaptersummary}


% ---------------------------------------------------------------------
% 以下は今後追加する（範囲表の大項目）
% \part{数学基礎}
% \part{計算基礎}
% \part{モデリング・AI と評価}

\backmatter
\printindex

\end{document}
